%% ============================================================
%% sections/a10-transparency.tex
%% H. R. 9510 Bill v5.0 - APPENDIX E. DEVELOPMENT TRANSPARENCY AND COMMIT
%% SCHEDULE. How Bill v5.0 was built.
%% Rendered visuals: Table 14 (bill version lineage v1.0-v5.0), Table 15 (the
%% nine-milestone commit schedule), Figure 11 (Mermaid build process, TikZ).
%% ============================================================
\phantomsection\label{toc:a10}
\addcontentsline{toc}{section}{Appendix E. Development Transparency and Commit Schedule}
\usccrosshead{APPENDIX E. DEVELOPMENT TRANSPARENCY AND COMMIT SCHEDULE}

\uscprov{1}{\textbf{What was built, and how.} Bill v5.0 (repository release
v2.0.0) is the Financial Data Amendment of H. R. 9510. It was produced by an
autonomous artificial intelligence coding agent (Claude Code, 1M context)
operating under human direction in a managed, ephemeral cloud container, from
a single Master prompt: Process A generated seven sub-prompts, and Process B
ran them in sequence through four research stages and three bill stages,
building from the Bill v3.0 measure (without its deliverables) and the
financial frame identified in Stage 1. It was authored across sequential
commits, one file per commit pushed to GitHub in real time, with the nine
stage milestones tracked in one continuously updated pull request on the
single permitted branch. Only kevinkawchak/single-prompt-bill was edited
\cite{anthropic-claude-code, repo-single-prompt-bill}.}

{\footnotesize
\begin{xltabular}{\textwidth}{@{}L{1.4cm} L{4.4cm} Y L{2.7cm}@{}}
\hline\noalign{\vskip 0.04cm}
\textbf{Version} & \textbf{Directory} & \textbf{Role} & \textbf{DOI}\\
\hline\noalign{\vskip 0.04cm}
\endfirsthead
\hline\noalign{\vskip 0.04cm}
\textbf{Version} & \textbf{Directory} & \textbf{Role} & \textbf{DOI}\\
\hline\noalign{\vskip 0.04cm}
\endhead
\hline\noalign{\vskip 0.04cm}
\endlastfoot
Bill v1.0 & cancer-automated: papers/VVUQ-03/final-paper & Standalone Physical
AI Oncology Trial Bill & \tiny 10.5281/zenodo.20454870\\
\hline\noalign{\vskip 0.04cm}
Bill v2.0 & cancer-automated: papers/VVUQ-04/final-bill & FD\&C Act amendment:
new section 360e-5 & \tiny 10.5281/zenodo.20485580\\
\hline\noalign{\vskip 0.04cm}
Bill v3.0 & cancer-automated: papers/VVUQ-05/final-bill & The visual amendment
this version is built from & \tiny 10.5281/zenodo.20535429\\
\hline\noalign{\vskip 0.04cm}
Bill v4.0 & single-prompt-bill: auto-bill-01 & The gray-scale Mermaid
amendment; the process adapted here & \tiny 10.5281/zenodo.20576907\\
\hline\noalign{\vskip 0.04cm}
Bill v5.0 & single-prompt-bill: auto-bill-02 & This Financial Data Amendment &
\tiny 10.5281/zenodo.xxxxxxxx\\
\hline
\end{xltabular}}
\vspace{-1cm}
\figcaption{Table 14. Bill version lineage: each bill version, its directory,
its role, and its DOI (the v5.0 DOI is the Rule 15 placeholder pending
deposit).}

\uscprov{1}{\textbf{What this version adds.} From Bill v3.0, this version
keeps the full operative text and adds the financial perspective: the section
515D(k) financial-data record, the SEC. 5 financial core (transparency, the
user-fee treatment, the authorization of appropriations, the boxed
budgetary-effects clause), the twelfth comparative-print section, and the
four financial appendices, rendered in the three media - fifteen full-width
tables, six ASCII figures with the cover, and six gray-scale Mermaid diagrams
as TikZ \cite{kawchak2026vvuq05bill, kawchak2026autobill01}. The ordered
milestones of the single pull request are listed in Table 15; Figure 11 draws
the build.}

{\footnotesize
\begin{xltabular}{\textwidth}{@{}L{2.2cm} L{5.6cm} Y@{}}
\hline\noalign{\vskip 0.04cm}
\textbf{Milestone} & \textbf{Files or action} & \textbf{Role}\\
\hline\noalign{\vskip 0.04cm}
\endfirsthead
\hline\noalign{\vskip 0.04cm}
\textbf{Milestone} & \textbf{Files or action} & \textbf{Role}\\
\hline\noalign{\vskip 0.04cm}
\endhead
\hline\noalign{\vskip 0.04cm}
\endlastfoot
M1 (10 commits) & README, master-prompt, seven sub-prompts; the PR opened &
Bootstrap; Process A\\
\hline\noalign{\vskip 0.04cm}
M2-M5 (2 each) & 01-research through 04-figure-selection, output plus README &
The four research stages\\
\hline\noalign{\vskip 0.04cm}
M6 (16 commits) & draft-bill: main, style, bib, nine scaffolds, prompt,
output, README, error-pass plus zip & Stage 5, the scaffold\\
\hline\noalign{\vskip 0.04cm}
M7 (16 commits) & full-bill: the same set, every slot rendered & Stage 6, this
full bill\\
\hline\noalign{\vskip 0.04cm}
M8 (16 commits) & final-bill: the same set, with the senior-author polish & Stage 7; repository v2.0.0\\
\hline\noalign{\vskip 0.04cm}
M9 (4 commits) & CHANGELOG.md, releases.md, root README, build status & The
single last release update\\
\hline
\end{xltabular}}
\vspace{-0.8cm}
\figcaption{Table 15. The nine-milestone commit schedule of this build: one
commit per file pushed in real time, every milestone a contiguous labeled set
inside the single continuously updated pull request.}

\begin{mermaidfig}
\node[mmgoal,text width=30mm] (mp) at (0,0) {Master prompt};
\node[mmstep,text width=42mm] (sp) at (0,-1.8) {M1: seven sub-prompts,\\READMEs, the single PR};
\node[mmin,text width=26mm]  (r)  at (-4.4,-3.7) {M2 research};
\node[mmin,text width=26mm]  (t)  at (0,-3.7)    {M3 template};
\node[mmin,text width=26mm]  (ma) at (4.4,-3.7)  {M4 Mermaid + ASCII};
\node[mmin,text width=26mm]  (f)  at (-4.4,-5.5) {M5 figures};
\node[mmstep,text width=26mm] (d)  at (0,-5.5)   {M6 draft bill};
\node[mmlaw,text width=26mm] (fu) at (4.4,-5.5)  {M7 full bill};
\node[mmlaw,text width=26mm] (fi) at (0,-7.6)    {M8 final bill};
\node[mmgoal,text width=34mm](rel) at (4.4,-7.6) {M9 release v2.0.0};
\draw[mmedge] (mp) -- (sp);
\draw[mmedge] (sp) -- (r);
\draw[mmedge] (r) -- (t);
\draw[mmedge] (t) -- (ma);
\draw[mmedge] (ma.south) -- ++(0,-0.45) -- ++(-8.8,0) -- (f.north);
\draw[mmedge] (f) -- (d);
\draw[mmedge] (d) -- (fu);
\draw[mmedge] (fu) -- (fi);
\draw[mmedge] (fi) -- (rel);
\end{mermaidfig}
\vspace{-0.2cm}
\figcaption{Figure 11. The nine-milestone build process (gray-scale Mermaid as
TikZ): Process A at the top, the Process B stages in sequence, and the release
milestone that closes the repository at v2.0.0.}

\uscprov{1}{\textbf{How this measure is implemented as law.} On enactment, the
Secretary issues final regulations within 365 days under section 515D(h); the
amendments take effect on the first day of the first calendar quarter
beginning not less than one year after enactment; an investigation already
enrolling has 180 days to come into compliance; the first annual
financial-data report under section 515D(k)(3) is due for the first full
fiscal year after the effective date, and the first de-identified public
summary follows under section 5(a); the authorization of appropriations runs
fiscal years 2027 through 2031 on the Table 6 schedule, subject to
appropriations Acts; and any verification-review workload is priced in the
MDUFA VI negotiation due before October 1, 2027. The enforcement trigger
remains a single fact - the absence of a cleared
verification-before-generation record - and the financial trigger is now its
twin: the absence of the cost record that section 515D(k) attaches to it. The
financial-data standard of Appendix C grades this appendix's own commit
record, so the transparency the bill demands of sponsors is the transparency
this build practiced \cite{pl-mdufa5, gao-greenbook}.}
